% filepath: /home/kluo/Documents/Repos/mathphy/homework/2026/hw1_sol.tex
% Solutions to Homework 1 — Math Physics Methods (数学物理方法)

\PassOptionsToPackage{quiet}{xeCJK}
\documentclass[11pt]{article}

\usepackage[margin=1in]{geometry}
\usepackage{amsmath,amsthm,amssymb,mathtools, graphicx, multicol, array}
\usepackage{ctex}
\usepackage{fancyhdr}
\usepackage{fontspec}
\usepackage{fourier-orns}
\usepackage{tikz}
\usetikzlibrary{decorations.markings}

% Chinese punctuation handling (same as hw1.tex)
\let\oldjuhao=。
\catcode`\。=\active
\newcommand{。}{\ifmmode\text{\oldjuhao}\else\oldjuhao\fi}

\let\olddouhao=，
\catcode`\，=\active
\newcommand{，}{\ifmmode\text{\olddouhao}\else\olddouhao\fi}

% Common math symbols
\newcommand{\N}{\mathbb{N}}
\newcommand{\Z}{\mathbb{Z}}
\newcommand{\R}{\mathbb{R}}
\newcommand{\C}{\mathbb{C}}

% Theorem environments
\theoremstyle{remark}
\newtheorem{problem}{}
\theoremstyle{definition}
\newtheorem{solution}{解}

\renewcommand{\headrule}{%
\vspace{-8pt}\hrulefill
\raisebox{-2.0pt}{\quad\decofourleft\decotwo\decofourright\quad}\hrulefill}

% Silence fancyhdr headheight warning (same as hw1.tex)
\setlength{\headheight}{15.11996pt}
\setlength{\topmargin}{-3.11996pt}

\newcommand{\duedate}{{\zhdate{2026/9/18}}}
\pagestyle{fancy}

\fancyhead{} % clear all header fields
\fancyhead[L]{数学物理方法}
\zhnumsetup{style={Traditional,Financial}}
\fancyhead[R]{课后习题\zhnumber{1} \quad 参考解答}

\fancyfoot{} % clear all footer fields
\fancyfoot[C]{\thepage}
\fancyfoot[L]{截止日期\duedate}
\fancyfoot[R]{任课教师:罗凯}

\begin{document}
\renewcommand{\labelenumi}{(\arabic{enumi})}
\renewcommand{\labelenumii}{(\arabic{enumi}.\arabic{enumii})}

%==========================================================================
\begin{problem}
写出下列复数的实部、虚部、模和辐角。
\begin{enumerate}
    \item $1 + \imath \sqrt{3}$,
    \item $\dfrac{1-\imath}{1+\imath}$,
    \item $e^{z}$（其中 $z = x + \imath y$）,
    \item $e^{1+\imath}$,
    \item $e^{\imath \varphi(x)}$，其中 $\varphi(x)$ 是实数 $x$ 的实函数。
\end{enumerate}
\end{problem}

\begin{solution}
对每个复数 $z$，设其实部为 $\mathrm{Re}\,z$、虚部为 $\mathrm{Im}\,z$、模为 $|z|$、辐角为 $\mathrm{Arg}\,z$（取主值）。

\textbf{(1)} $z_1 = 1 + \imath\sqrt{3}$.\\
$\mathrm{Re}\,z_1 = 1$，\quad $\mathrm{Im}\,z_1 = \sqrt{3}$，\quad
$|z_1| = \sqrt{1^2 + (\sqrt{3})^2} = 2$，\quad
$\mathrm{Arg}\,z_1 = \arctan\!\big(\sqrt{3}/1\big) = \dfrac{\pi}{3}$.

\textbf{(2)} $z_2 = \dfrac{1-\imath}{1+\imath}$. 分子分母同乘 $(1-\imath)$：
\[
z_2 = \frac{(1-\imath)^2}{1^2 + 1^2} = \frac{1 - 2\imath + \imath^2}{2} = \frac{-2\imath}{2} = -\imath .
\]
故 $\mathrm{Re}\,z_2 = 0$，\quad $\mathrm{Im}\,z_2 = -1$，\quad $|z_2| = 1$，\quad
$\mathrm{Arg}\,z_2 = -\dfrac{\pi}{2}$（亦可写为 $\dfrac{3\pi}{2}$）.

\textbf{(3)} $z_3 = e^{z} = e^{x + \imath y} = e^x(\cos y + \imath \sin y)$.\\
$\mathrm{Re}\,z_3 = e^x \cos y$，\quad $\mathrm{Im}\,z_3 = e^x \sin y$，\quad
$|z_3| = e^x$，\quad $\mathrm{Arg}\,z_3 = y + 2k\pi$，$k \in \mathbb{Z}$.

\textbf{(4)} $z_4 = e^{1+\imath} = e \cdot e^{\imath} = e(\cos 1 + \imath \sin 1)$.\\
$\mathrm{Re}\,z_4 = e\cos 1$，\quad $\mathrm{Im}\,z_4 = e\sin 1$，\quad $|z_4| = e$，\quad $\mathrm{Arg}\,z_4 = 1$（弧度）.

\textbf{(5)} $z_5 = e^{\imath \varphi(x)} = \cos\varphi(x) + \imath \sin\varphi(x)$.\\
$\mathrm{Re}\,z_5 = \cos\varphi(x)$，\quad $\mathrm{Im}\,z_5 = \sin\varphi(x)$，\quad $|z_5| = 1$，\quad $\mathrm{Arg}\,z_5 = \varphi(x) + 2k\pi$，$k \in \mathbb{Z}$.
\end{solution}

%==========================================================================
\begin{problem}
证明以下规律：
\begin{enumerate}
    \item 复数的加法和乘法满足结合律，即
    \begin{itemize}
        \item $(z_1 + z_2) + z_3 = z_1 + (z_2 + z_3)$，
        \item $(z_1 \cdot z_2) \cdot z_3 = z_1 \cdot (z_2 \cdot z_3)$；
    \end{itemize}
    \item 复数乘积的共轭等于复数共轭的乘积，即 $(z_1 \cdot z_2)^* = z_1^*\, z_2^*$.
\end{enumerate}
\end{problem}

\begin{solution}
设 $z_k = x_k + \imath y_k$（$k=1,2,3$），则 $z_k^* = x_k - \imath y_k$.

\textbf{加法结合律：}
\begin{align*}
(z_1 + z_2) + z_3
&= \bigl[(x_1 + x_2) + \imath(y_1 + y_2)\bigr] + (x_3 + \imath y_3) \\
&= (x_1 + x_2 + x_3) + \imath(y_1 + y_2 + y_3), \\[4pt]
z_1 + (z_2 + z_3)
&= (x_1 + \imath y_1) + \bigl[(x_2 + x_3) + \imath(y_2 + y_3)\bigr] \\
&= (x_1 + x_2 + x_3) + \imath(y_1 + y_2 + y_3).
\end{align*}
两边相等，故 $(z_1 + z_2) + z_3 = z_1 + (z_2 + z_3)$.

\textbf{乘法结合律：}
先计算 $z_1 z_2$ 与 $z_2 z_3$：
\begin{align*}
z_1 z_2 &= (x_1 x_2 - y_1 y_2) + \imath(x_1 y_2 + x_2 y_1), \\
z_2 z_3 &= (x_2 x_3 - y_2 y_3) + \imath(x_2 y_3 + x_3 y_2).
\end{align*}
再分别与剩余的因子相乘：
\begin{align*}
(z_1 z_2) z_3 &=
\bigl[(x_1 x_2 - y_1 y_2) + \imath(x_1 y_2 + x_2 y_1)\bigr](x_3 + \imath y_3) \\
&= (x_1 x_2 - y_1 y_2)x_3 - (x_1 y_2 + x_2 y_1) y_3
   + \imath\bigl[(x_1 y_2 + x_2 y_1)x_3 + (x_1 x_2 - y_1 y_2) y_3\bigr] \\
&= (x_1 x_2 x_3 - y_1 y_2 x_3 - x_1 y_2 y_3 - x_2 y_1 y_3) \\
&\quad + \imath(x_1 y_2 x_3 + x_2 y_1 x_3 + x_1 x_2 y_3 - y_1 y_2 y_3).
\end{align*}
\begin{align*}
z_1 (z_2 z_3) &=
(x_1 + \imath y_1)\bigl[(x_2 x_3 - y_2 y_3) + \imath(x_2 y_3 + x_3 y_2)\bigr] \\
&= x_1(x_2 x_3 - y_2 y_3) - y_1(x_2 y_3 + x_3 y_2)
   + \imath\bigl[y_1(x_2 x_3 - y_2 y_3) + x_1(x_2 y_3 + x_3 y_2)\bigr] \\
&= (x_1 x_2 x_3 - x_1 y_2 y_3 - x_2 y_1 y_3 - x_3 y_1 y_2) \\
&\quad + \imath(x_1 x_2 y_3 + x_1 x_3 y_2 + x_2 x_3 y_1 - y_1 y_2 y_3).
\end{align*}
实部与虚部分别相等，故 $(z_1 z_2) z_3 = z_1 (z_2 z_3)$.

\textbf{共轭与乘法的交换性：}
\begin{align*}
(z_1 z_2)^*
   &= \bigl[(x_1 x_2 - y_1 y_2) + \imath(x_1 y_2 + x_2 y_1)\bigr]^* \\
   &= (x_1 x_2 - y_1 y_2) - \imath(x_1 y_2 + x_2 y_1), \\[4pt]
z_1^* z_2^*
   &= (x_1 - \imath y_1)(x_2 - \imath y_2) \\
   &= (x_1 x_2 - y_1 y_2) - \imath(x_1 y_2 + x_2 y_1).
\end{align*}
两边相等，故 $(z_1 z_2)^* = z_1^* z_2^*$. \hfill $\square$
\end{solution}

%==========================================================================
\begin{problem}
找出复平面上满足方程
\[
|z - \imath a| = \lambda\,|z + \imath a|,\qquad \lambda > 0
\]
的所有点 $z = x + \imath y$。请分 $\lambda$ 的三种情况
(1) $\lambda < 1$，\quad (2) $\lambda > 1$，\quad (3) $\lambda = 1$
分别绘制图形。
\end{problem}

\begin{solution}
令 $z = x + \imath y$，则
\[
|z - \imath a|^2 = x^2 + (y - a)^2,\qquad |z + \imath a|^2 = x^2 + (y + a)^2.
\]
原方程取平方得
\[
x^2 + (y - a)^2 = \lambda^2\bigl[x^2 + (y + a)^2\bigr].
\]
展开整理：
\begin{align*}
(1 - \lambda^2)(x^2 + y^2) - 2a(1 + \lambda^2)y + (1 - \lambda^2)a^2 &= 0.
\end{align*}

\textbf{情形 1：$\lambda < 1$.} \quad 此时 $1 - \lambda^2 > 0$，上式两边除以 $(1 - \lambda^2)$：
\[
x^2 + y^2 - \frac{2a(1 + \lambda^2)}{1 - \lambda^2}\,y + a^2 = 0,
\]
配方得
\[
x^2 + \Bigl(y - \frac{a(1 + \lambda^2)}{1 - \lambda^2}\Bigr)^2
   = \Bigl(\frac{2a\lambda}{1 - \lambda^2}\Bigr)^2.
\]
这是\textbf{圆心} $O_1 = \bigl(0,\, \frac{a(1+\lambda^2)}{1-\lambda^2}\bigr)$、
\textbf{半径} $r_1 = \dfrac{2a\lambda}{1-\lambda^2}$ 的圆，
位于实轴上方（因 $\lambda<1$ 使圆心 $y$ 坐标为正）。

\textbf{情形 2：$\lambda > 1$.} \quad 此时 $1 - \lambda^2 < 0$，除以 $1 - \lambda^2$ 后得
\[
x^2 + y^2 + \frac{2a(1 + \lambda^2)}{\lambda^2 - 1}\,y - a^2 = 0,
\]
配方得
\[
x^2 + \Bigl(y + \frac{a(1 + \lambda^2)}{\lambda^2 - 1}\Bigr)^2
   = \Bigl(\frac{2a\lambda}{\lambda^2 - 1}\Bigr)^2.
\]
圆心 $O_2 = \bigl(0,\, -\frac{a(1+\lambda^2)}{\lambda^2-1}\bigr)$，半径 $r_2 = \dfrac{2a\lambda}{\lambda^2-1}$，位于实轴下方。

\textbf{情形 3：$\lambda = 1$.} \quad 原方程化为 $|z - \imath a| = |z + \imath a|$，
表示与两点 $\imath a$ 和 $-\imath a$ 等距的点的轨迹，是\textbf{实轴}（即直线 $y = 0$）。

\textbf{几何说明。}
方程刻画的是阿波罗尼乌斯圆 (Apollonius circle)：平面上到两定点 $A=(0,a)$ 与
$B=(0,-a)$ 距离之比为 $\lambda$ 的点的轨迹。
当 $\lambda < 1$ 时圆在 $A$ 侧（更靠近 $A$）；当 $\lambda > 1$ 时圆在 $B$ 侧；当 $\lambda = 1$ 时退化为 $AB$ 的中垂线。
当 $\lambda \to 0$ 时圆心趋于 $A$，半径趋于 $0$；
当 $\lambda \to \infty$ 时圆心趋于 $B$，半径趋于 $0$。
\end{solution}

\begin{center}
% Reciprocal ratios \lambda=1/2 and \lambda=2 give identical radii 4/3 (a=1).
% Same bounding box so the two circles appear the same size.
\begin{tikzpicture}[scale=0.85, baseline=(current bounding box.center)]
  \useasboundingbox (-2.25,-3.65) rectangle (2.45,3.65);
  \draw[->,gray] (-2.1,0) -- (2.2,0) node[right,black] {$\mathrm{Re}$};
  \draw[->,gray] (0,-3.35) -- (0,3.45) node[above,black] {$\mathrm{Im}$};
  \draw[thick,blue!70!black] (0,5/3) circle[radius=4/3];
  \fill (0,5/3) circle[radius=1.4pt];
  \node[anchor=west] at (0.22,5/3) {$O_1$};
  \fill (0,1) circle[radius=1.4pt];
  \node[anchor=east] at (-0.22,1) {$\imath a$};
  \fill (0,-1) circle[radius=1.4pt];
  \node[anchor=east] at (-0.22,-1) {$-\imath a$};
  \node at (0,-3.52) {$\lambda=1/2<1$};
\end{tikzpicture}
\hspace{0.9cm}
\begin{tikzpicture}[scale=0.85, baseline=(current bounding box.center)]
  \useasboundingbox (-2.25,-3.65) rectangle (2.45,3.65);
  \draw[->,gray] (-2.1,0) -- (2.2,0) node[right,black] {$\mathrm{Re}$};
  \draw[->,gray] (0,-3.35) -- (0,3.45) node[above,black] {$\mathrm{Im}$};
  \draw[thick,blue!70!black] (0,-5/3) circle[radius=4/3];
  \fill (0,-5/3) circle[radius=1.4pt];
  \node[anchor=west] at (0.22,-5/3) {$O_2$};
  \fill (0,1) circle[radius=1.4pt];
  \node[anchor=east] at (-0.22,1) {$\imath a$};
  \fill (0,-1) circle[radius=1.4pt];
  \node[anchor=east] at (-0.22,-1) {$-\imath a$};
  \node at (0,-3.52) {$\lambda=2>1$};
\end{tikzpicture}
\hspace{0.9cm}
\begin{tikzpicture}[scale=0.85, baseline=(current bounding box.center)]
  \useasboundingbox (-2.25,-3.65) rectangle (2.45,3.65);
  \draw[->,gray] (-2.1,0) -- (2.2,0) node[right,black] {$\mathrm{Re}$};
  \draw[->,gray] (0,-3.35) -- (0,3.45) node[above,black] {$\mathrm{Im}$};
  \draw[very thick,blue!70!black] (-1.85,0) -- (1.85,0);
  \fill (0,1) circle[radius=1.4pt];
  \node[anchor=east] at (-0.22,1) {$\imath a$};
  \fill (0,-1) circle[radius=1.4pt];
  \node[anchor=east] at (-0.22,-1) {$-\imath a$};
  \node at (0,-3.52) {$\lambda=1$};
\end{tikzpicture}
\end{center}

%==========================================================================
\begin{problem}
若 $|z| = 1$，$a, b$ 为任意复数，试证明
\[
\left| \frac{a z + b}{\bar{b} z + \bar{a}} \right| = 1.
\]
\end{problem}

\begin{solution}
只需证明 $|a z + b|^2 = |\bar{b} z + \bar{a}|^2$.

\begin{align*}
|a z + b|^2
&= (a z + b)\,\overline{(a z + b)}
 = (a z + b)(\bar{a}\,\bar{z} + \bar{b}) \\
&= a\bar{a}\,z\bar{z} + a\bar{b}\,z + b\bar{a}\,\bar{z} + b\bar{b} \\
&= |a|^2 + |b|^2 + a\bar{b}\,z + b\bar{a}\,\bar{z} \quad (\text{因为 } |z|=1 \Rightarrow z\bar z = 1).
\end{align*}

\begin{align*}
|\bar{b} z + \bar{a}|^2
&= (\bar{b} z + \bar{a})\,\overline{(\bar{b} z + \bar{a})}
 = (\bar{b} z + \bar{a})(b\bar{z} + a) \\
&= \bar{b} b\,z\bar{z} + \bar{b} a\,z + \bar{a} b\,\bar{z} + \bar{a} a \\
&= |a|^2 + |b|^2 + \bar{b} a\,z + \bar{a} b\,\bar{z}.
\end{align*}

由于复数乘法可交换，$\bar{b} a = a\bar{b}$，$\bar{a} b = b\bar{a}$，故
\[
a\bar{b}\,z + b\bar{a}\,\bar{z} = \bar{b} a\,z + \bar{a} b\,\bar{z}.
\]
从而
\[
|a z + b|^2 = |\bar{b} z + \bar{a}|^2
\;\Longrightarrow\;
|a z + b| = |\bar{b} z + \bar{a}|.
\]
当 $\bar{b} z + \bar{a} \neq 0$ 时，两边同除之即得
$\left|\dfrac{a z + b}{\bar{b} z + \bar{a}}\right| = 1$. \hfill $\square$
\end{solution}

%==========================================================================
\begin{problem}
试用 $\cos\varphi,\ \sin\varphi$ 表示 $\cos 4\varphi$.
\end{problem}

\begin{solution}
\textbf{方法一（直接用二倍角公式）。}
\begin{align*}
\cos 4\varphi
&= \cos^2 2\varphi - \sin^2 2\varphi \\
&= (2\cos^2\varphi - 1)^2 - (2\sin\varphi\cos\varphi)^2 \\
&= 4\cos^4\varphi - 4\cos^2\varphi + 1 - 4\sin^2\varphi\cos^2\varphi \\
&= 4\cos^4\varphi - 4\cos^2\varphi + 1 - 4\cos^2\varphi(1 - \cos^2\varphi) \\
&= 8\cos^4\varphi - 8\cos^2\varphi + 1.
\end{align*}

\textbf{方法二（用 $e^{\imath\varphi}$）。}
\begin{align*}
\cos 4\varphi &= \mathrm{Re}\!\bigl(e^{\imath 4\varphi}\bigr)
              = \mathrm{Re}\!\bigl((\cos\varphi + \imath\sin\varphi)^4\bigr) \\
             &= \cos^4\varphi - 6\cos^2\varphi\sin^2\varphi + \sin^4\varphi \\
             &= (\cos^2\varphi + \sin^2\varphi)^2 - 8\cos^2\varphi\sin^2\varphi \\
             &= 1 - 2\sin^2 2\varphi \\
             &= 1 - 8\sin^2\varphi\cos^2\varphi.
\end{align*}

\textbf{两种表达式的等价性。}
\[
1 - 8\sin^2\varphi\cos^2\varphi
= 1 - 2(2\sin\varphi\cos\varphi)^2
= 1 - 2\sin^2 2\varphi
= \cos 4\varphi,
\]
又
\[
8\cos^4\varphi - 8\cos^2\varphi + 1
= 1 - 8\cos^2\varphi(1 - \cos^2\varphi)
= 1 - 8\sin^2\varphi\cos^2\varphi.
\]

\[
\boxed{\cos 4\varphi = 8\cos^4\varphi - 8\cos^2\varphi + 1
       = 1 - 8\sin^2\varphi\cos^2\varphi
       = 1 - 2\sin^2 2\varphi.}
\]
\end{solution}

%==========================================================================
\begin{problem}
求下列方程的根，并在复平面上画出它们的位置。
\begin{enumerate}
    \item $z^4 + 1 = 0$；
    \item $z^2 + 2z\cos\lambda + 1 = 0$，\quad $0 < \lambda < \pi$.
\end{enumerate}
\end{problem}

\begin{solution}

\textbf{(1) $z^4 + 1 = 0$.}
即 $z^4 = -1 = e^{\imath(\pi + 2k\pi)}$，$k = 0, 1, 2, 3$，故
\[
z_k = e^{\imath(\pi + 2k\pi)/4} = e^{\imath\pi(2k+1)/4},\qquad k=0,1,2,3.
\]
显式列出：
\begin{align*}
z_0 &= e^{\imath\pi/4} = \tfrac{\sqrt 2}{2}(1 + \imath), &
z_1 &= e^{\imath 3\pi/4} = \tfrac{\sqrt 2}{2}(-1 + \imath), \\
z_2 &= e^{\imath 5\pi/4} = \tfrac{\sqrt 2}{2}(-1 - \imath), &
z_3 &= e^{\imath 7\pi/4} = \tfrac{\sqrt 2}{2}(1 - \imath).
\end{align*}
四点均位于单位圆上，彼此相隔 $\pi/2$（在第一、二、三、四象限各一个）。

\textbf{(2) $z^2 + 2z\cos\lambda + 1 = 0$，$0 < \lambda < \pi$.}
由求根公式：
\begin{align*}
z &= \frac{-2\cos\lambda \pm \sqrt{4\cos^2\lambda - 4}}{2}
   = -\cos\lambda \pm \sqrt{\cos^2\lambda - 1} \\
  &= -\cos\lambda \pm \sqrt{-\sin^2\lambda}
   = -\cos\lambda \pm \imath\sin\lambda.
\end{align*}
故
\[
z_1 = -\cos\lambda + \imath\sin\lambda,\qquad
z_2 = -\cos\lambda - \imath\sin\lambda.
\]
模与辐角：
$|z_1| = |z_2| = \sqrt{\cos^2\lambda + \sin^2\lambda} = 1$，两根都在单位圆上。
记 $\theta = \pi - \lambda$（因 $0<\lambda<\pi$ 故 $0<\theta<\pi$），则
\[
z_1 = e^{\imath(\pi-\lambda)},\qquad z_2 = e^{-\imath(\pi-\lambda)} = e^{\imath(\lambda - \pi)},
\]
即两根关于实轴对称，位于单位圆上辐角分别为 $\pm(\pi-\lambda)$ 的两点。

\begin{center}
\begin{tikzpicture}[scale=1.35, baseline=(current bounding box.center)]
  \useasboundingbox (-1.75,-1.95) rectangle (1.75,1.75);
  \draw[->,gray] (-1.55,0) -- (1.55,0) node[right,black] {$\mathrm{Re}$};
  \draw[->,gray] (0,-1.55) -- (0,1.55) node[above,black] {$\mathrm{Im}$};
  \draw (0,0) circle[radius=1];
  \foreach \ang/\lab in {45/$z_0$, 135/$z_1$, 225/$z_2$, 315/$z_3$} {
    \fill (\ang:1) circle[radius=1.2pt];
    \node at (\ang:1.32) {\lab};
  }
  \node at (0,-1.82) {\small $z^4+1=0$};
\end{tikzpicture}
\hspace{1.1cm}
\begin{tikzpicture}[scale=1.35, baseline=(current bounding box.center)]
  \useasboundingbox (-1.75,-1.95) rectangle (1.75,1.75);
  \draw[->,gray] (-1.55,0) -- (1.55,0) node[right,black] {$\mathrm{Re}$};
  \draw[->,gray] (0,-1.55) -- (0,1.55) node[above,black] {$\mathrm{Im}$};
  \draw (0,0) circle[radius=1];
  \draw[dashed] (0,0) -- (120:1);
  \draw[dashed] (0,0) -- (-120:1);
  \fill (120:1) circle[radius=1.2pt];
  \node at (150:1.32) {$z_1$};
  \fill (-120:1) circle[radius=1.2pt];
  \node at (-150:1.32) {$z_2$};
  \node at (0,-1.82) {\small $\lambda=\pi/3$};
\end{tikzpicture}
\end{center}
\end{solution}

%==========================================================================
\begin{problem}
验证
\[
\tan^{-1}(z) = \frac{\imath}{2}\bigl[\ln(1 - \imath z) - \ln(1 + \imath z)\bigr].
\]
\end{problem}

\begin{solution}
设 $w = \tan^{-1}(z)$，则 $z = \tan w = \dfrac{\sin w}{\cos w}$.
利用欧拉公式
\[
\sin w = \frac{e^{\imath w} - e^{-\imath w}}{2\imath},\qquad
\cos w = \frac{e^{\imath w} + e^{-\imath w}}{2},
\]
得
\[
z = \frac{e^{\imath w} - e^{-\imath w}}{\imath\bigl(e^{\imath w} + e^{-\imath w}\bigr)}.
\]
令 $u = e^{\imath w}$，则 $e^{-\imath w} = u^{-1}$，于是
\[
z = \frac{u - u^{-1}}{\imath(u + u^{-1})}
   = \frac{u^2 - 1}{\imath(u^2 + 1)}.
\]
整理求解 $u^2$：
\[
\imath z(u^2 + 1) = u^2 - 1
\;\Longrightarrow\;
u^2(\imath z - 1) = -1 - \imath z
\;\Longrightarrow\;
u^2 = \frac{1 + \imath z}{1 - \imath z}.
\]
因此（取主值）
\[
e^{\imath w} = u = \sqrt{\frac{1 + \imath z}{1 - \imath z}}
\;\Longrightarrow\;
\imath w = \frac{1}{2}\ln\frac{1 + \imath z}{1 - \imath z}.
\]
故
\[
w = \frac{1}{2\imath}\ln\frac{1 + \imath z}{1 - \imath z}
   = \frac{\imath}{2}\cdot \frac{1}{\imath^2}\ln\frac{1 + \imath z}{1 - \imath z}
   = -\frac{\imath}{2}\ln\frac{1 + \imath z}{1 - \imath z}.
\]
由对数差的性质 $\ln(A/B) = \ln A - \ln B$：
\[
w = -\frac{\imath}{2}\bigl[\ln(1 + \imath z) - \ln(1 - \imath z)\bigr]
   = \frac{\imath}{2}\bigl[\ln(1 - \imath z) - \ln(1 + \imath z)\bigr].
\]
即
\[
\boxed{\tan^{-1}(z) = \frac{\imath}{2}\bigl[\ln(1 - \imath z) - \ln(1 + \imath z)\bigr].}\qquad\square
\]
\end{solution}

%==========================================================================
\begin{problem}
试证明极坐标下的柯西-黎曼方程
\[
\begin{cases}
\dfrac{\partial u}{\partial \rho} = \dfrac{1}{\rho}\dfrac{\partial v}{\partial \varphi},\\[6pt]
\dfrac{1}{\rho}\dfrac{\partial u}{\partial \varphi} = -\dfrac{\partial v}{\partial \rho}.
\end{cases}
\]
\end{problem}

\begin{solution}
设 $z = \rho e^{\imath\varphi}$，则 $x = \rho\cos\varphi$，$y = \rho\sin\varphi$，故
\[
\frac{\partial x}{\partial \rho} = \cos\varphi,\quad
\frac{\partial x}{\partial \varphi} = -\rho\sin\varphi,\quad
\frac{\partial y}{\partial \rho} = \sin\varphi,\quad
\frac{\partial y}{\partial \varphi} = \rho\cos\varphi.
\]
由链式法则
\begin{align*}
\frac{\partial u}{\partial \rho}
   &= \frac{\partial u}{\partial x}\cos\varphi + \frac{\partial u}{\partial y}\sin\varphi, \\
\frac{\partial u}{\partial \varphi}
   &= -\rho\sin\varphi\,\frac{\partial u}{\partial x} + \rho\cos\varphi\,\frac{\partial u}{\partial y}, \\
\frac{\partial v}{\partial \rho}
   &= \frac{\partial v}{\partial x}\cos\varphi + \frac{\partial v}{\partial y}\sin\varphi, \\
\frac{\partial v}{\partial \varphi}
   &= -\rho\sin\varphi\,\frac{\partial v}{\partial x} + \rho\cos\varphi\,\frac{\partial v}{\partial y}.
\end{align*}
由直角坐标下的 Cauchy--Riemann 方程
$\dfrac{\partial u}{\partial x} = \dfrac{\partial v}{\partial y}$，
$\dfrac{\partial u}{\partial y} = -\dfrac{\partial v}{\partial x}$，
代入得
\[
\frac{1}{\rho}\frac{\partial v}{\partial \varphi}
   = -\sin\varphi\,\frac{\partial v}{\partial x} + \cos\varphi\,\frac{\partial v}{\partial y}
   = \sin\varphi\,\frac{\partial u}{\partial y} + \cos\varphi\,\frac{\partial u}{\partial x}
   = \frac{\partial u}{\partial \rho}.
\]
同理
\[
-\frac{\partial v}{\partial \rho}
   = -\cos\varphi\,\frac{\partial v}{\partial x} - \sin\varphi\,\frac{\partial v}{\partial y}
   = \cos\varphi\,\frac{\partial u}{\partial y} - \sin\varphi\,\frac{\partial u}{\partial x}.
\]
而
\[
\frac{1}{\rho}\frac{\partial u}{\partial \varphi}
   = -\sin\varphi\,\frac{\partial u}{\partial x} + \cos\varphi\,\frac{\partial u}{\partial y}.
\]
两式相等，故
\[
\frac{1}{\rho}\frac{\partial u}{\partial \varphi} = -\frac{\partial v}{\partial \rho}. \qquad\square
\]
\end{solution}

%==========================================================================
\begin{problem}
已知解析函数 $f(z)$ 的实部 $u(x, y)$ 或虚部 $v(x, y)$，求该解析函数。
\begin{enumerate}
    \item $u = e^x \sin y$；
    \item $u = x^2 - y^2 + xy$，\quad $f(0) = 0$.
\end{enumerate}
\end{problem}

\begin{solution}

\textbf{(1) $u = e^x \sin y$.}
\[
u_x = e^x \sin y,\qquad u_y = e^x \cos y.
\]
由 CR 方程 $v_y = u_x = e^x \sin y$，$v_x = -u_y = -e^x \cos y$.

对 $v_y$ 关于 $y$ 积分：
\[
v(x,y) = \int e^x \sin y\,\mathrm{d}y = -e^x \cos y + g(x).
\]
再由 $v_x = -u_y$：
\[
v_x = -e^x \cos y + g'(x) = -e^x \cos y \;\Longrightarrow\; g'(x) = 0 \;\Longrightarrow\; g(x) = C.
\]
所以
\[
v(x,y) = -e^x \cos y + C.
\]
合并：
\[
f(z) = u + \imath v
   = e^x \sin y + \imath\bigl(-e^x \cos y + C\bigr)
   = -\imath e^x(\cos y + \imath\sin y) + \imath C
   = -\imath e^{x + \imath y} + \imath C
   = -\imath e^z + \imath C.
\]
吸收常数后，$\boxed{f(z) = -\imath e^z}$（相差一个纯虚常数）。

\textbf{(2) $u = x^2 - y^2 + xy$，$f(0) = 0$.}
\[
u_x = 2x + y,\qquad u_y = -2y + x.
\]
由 CR 方程 $v_y = u_x = 2x + y$，$v_x = -u_y = 2y - x$.

对 $v_y$ 关于 $y$ 积分：
\[
v(x,y) = \int (2x + y)\,\mathrm{d}y = 2xy + \tfrac{1}{2}y^2 + g(x).
\]
再由 $v_x = -u_y$：
\[
v_x = 2y + g'(x) = 2y - x \;\Longrightarrow\; g'(x) = -x \;\Longrightarrow\; g(x) = -\tfrac{1}{2}x^2 + C.
\]
故
\[
v(x,y) = 2xy + \tfrac{1}{2}y^2 - \tfrac{1}{2}x^2 + C.
\]
合并：
\begin{align*}
f(z) &= u + \imath v \\
     &= (x^2 - y^2 + xy) + \imath\!\left(2xy + \tfrac{1}{2}y^2 - \tfrac{1}{2}x^2 + C\right) \\
     &= (x^2 - y^2) + \tfrac{\imath}{2}\bigl[(2\imath xy - \imath^2 x^2) + (2\imath xy + y^2)\bigr] \cdot 1 \; + \;\text{整理}\\
     &= \bigl[1 - \tfrac{\imath}{2}\bigr](x^2 - y^2 + 2\imath xy) + \imath C \\
     &= \Bigl(1 - \frac{\imath}{2}\Bigr)z^2 + \imath C.
\end{align*}
利用 $f(0) = 0$ 得 $\imath C = 0$，即 $C = 0$，从而
\[
\boxed{f(z) = \left(1 - \frac{\imath}{2}\right)z^2 = \frac{(2 - \imath)\,z^2}{2}.}
\]
\end{solution}

\end{document}